% Compile with: uplatex file.tex && dvipdfmx file.dvi
%           or: latexmk -uplatex file.tex
\documentclass[dvipdfmx]{article}
\usepackage{amsmath, amssymb}
\usepackage{tikz}

\title{scenario\_integration\_variable\_notation}
\date{\today}

\begin{document}
\maketitle

Register $\iint(x^2+y^2)\,dx\,dy$ over the unit disc $x^2+y^2\le1$. The CAS is not called yet — this is a range-bearing, display-only object.

definite-integrate $x^2$ + $y^2$ over the region $x^2$ + $y^2$ $\le$ 1

\begin{equation}
  \iint_{x^2+y^2 \le 1} (x^2 + y^2) \, dx \, dy
  \label{eq:E1}
\end{equation}

Change to polar coordinates, writing the angle as the Greek letter typed straight from the keyboard — that is also what the notebook displays, so being able to type it back is the natural thing. The Jacobian $|J|=r$ multiplies the integrand and the new bounds are derived automatically.

change of variables in E1 from x,y to r,θ via x=r*cos(θ), y=r*sin(θ)

\begin{equation}
  \int_{0}^{2\pi} \int_{0}^{1} r^3 \, dr \, d\theta
  \label{eq:E2}
\end{equation}

Evaluate the transformed integral.

evaluate this integral

\begin{equation}
  \int_{0}^{2\pi} \int_{0}^{1} r^3 \, dr \, d\theta = \frac{\pi}{2}
  \label{eq:E3}
\end{equation}

Now the same calculation in TeX spelling, to confirm that typing the displayed form back in gives the same answer.

definite-integrate $x^2$ + $y^2$ over the region $x^2$ + $y^2$ $\le$ 1

\begin{equation}
  \iint_{x^2+y^2 \le 1} (x^2 + y^2) \, dx \, dy
  \label{eq:E4}
\end{equation}

change of variables in E4 from x,y to r,$\theta$ via x=r*$\cos$($\theta$), y=r*$\sin$($\theta$)

\begin{equation}
  \int_{0}^{2\pi} \int_{0}^{1} r^3 \, dr \, d\theta
  \label{eq:E5}
\end{equation}

evaluate this integral

\begin{equation}
  \int_{0}^{2\pi} \int_{0}^{1} r^3 \, dr \, d\theta = \frac{\pi}{2}
  \label{eq:E6}
\end{equation}

The same for a single-variable definite integral, with both the variable and the bound written as Greek letters.

$sin(\theta )^2$

\begin{equation}
  \sin(\theta)^2
  \label{eq:E7}
\end{equation}

integrate E7 with respect to θ from 0 to π

\begin{equation}
  \int_{0}^{\pi} \sin(\theta)^2 \, d\theta = \frac{\pi}{2}
  \label{eq:E8}
\end{equation}

\end{document}